\section{Harapan Hasil Peramalan}
Hasil akhir dari proses pemodelan STARMA divisualisasikan dalam bentuk grafik yang membandingkan nilai prediksi dengan data aktual curah hujan. Visualisasi dilakukan pada lima titik observasi yang mewakili bagian-bagian wilayah Surabaya, yaitu Surabaya Barat, Timur, Utara, Selatan, dan Pusat. Hasil prediksi ini menjadi dasar untuk mengevaluasi sejauh mana model mampu merepresentasikan dinamika spasio-temporal curah hujan di Surabaya secara akurat.